% Generated by build_appendix.py: the authors' prompts, verbatim
\promptentry{1}{Setup}{2026-10-22 08:34 UTC}{Hey Claude, take a look at the project and the INSTRUCTIONS.md file and confirm that you understand how we will proceed.\newline If any questions, ask them}
\promptentry{2}{Setup}{2026-10-22 08:40 UTC}{1. No environment for now, you can build one.\newline 2. Yes\newline 3. perfect\newline 4. Indeed but let's not sacrifice quality}
\promptentry{3}{Step 1}{2026-10-22 08:51 UTC}{Go}
\promptentry{4}{Step 1}{2026-10-22 09:04 UTC}{Can we add a correlation matrix to the EDA ?}
\promptentry{5}{Step 1}{2026-10-22 09:09 UTC}{One last detail, the markdown rendering for mathematical expression in 1.1 are not great (example : C\_FN etc), do you think you could apply actual adapted mathematical markdown for these expressions'}
\promptentry{6}{Step 2}{2026-10-22 09:11 UTC}{Excellent, let's move on to part 2.}
\promptentry{7}{Step 2}{2026-10-22 09:30 UTC}{Would it be sound to do as in the course (lowest cross-validated log-loss) and compare it with the method of tuning it together with the tree in step 3 ?}
\promptentry{8}{Step 2}{2026-10-22 09:31 UTC}{Go}
\promptentry{9}{Step 2}{2026-10-22 09:44 UTC}{Can you explain the bootstrap selection frequency plot ? Im not sure to understand it. \newline Explain the two colors}
\promptentry{10}{Step 2}{2026-10-22 09:54 UTC}{Yes make that change.}
\promptentry{11}{Step 3}{2026-10-22 09:57 UTC}{Excellent, let's move on to the step 3 plan.}
\promptentry{12}{Step 3}{2026-10-22 10:05 UTC}{Go, let's follow your recommendation}
\promptentry{13}{Step 4}{2026-10-22 11:21 UTC}{Let's move on to the step 4 plan}
\promptentry{14}{Step 4}{2026-10-22 11:28 UTC}{1. Only the final model and the baseline\newline 2. explain what you mean \newline 3. explain what you mean}
\promptentry{15}{Step 4}{2026-10-22 11:30 UTC}{Go}
\promptentry{16}{Final review}{2026-10-22 12:04 UTC}{Great, now I want you to take a step back, review the full codebase, then the notebook. Do not modify a single thing, I just want you to look at it, be critical, and come up with a fair judgement of what we have done. Are there an pothholes, is it solid and rigorous ? \newline Evaluate and come back here in the discussion with a report. \newline Be concize and straight to the point. \newline If any issues, suggest improvements.}
\promptentry{17}{Final review}{2026-10-22 12:18 UTC}{For point 3, is the absent evidence important or can we simply get rid of it ?}
\promptentry{18}{Final review}{2026-10-22 12:19 UTC}{As you know, our goal is to make the work the simplest possible, therefore, get rid of everything you can get rid of without destroying the reasoning.}
\promptentry{19}{Final review}{2026-10-22 12:23 UTC}{What could we expect from fixing number 2 ?}
\promptentry{20}{Final review}{2026-10-22 12:24 UTC}{Go for it}
\promptentry{21}{Final review}{2026-10-22 12:24 UTC}{what would fixing \#1 look like ?}
\promptentry{22}{Final review}{2026-10-22 12:25 UTC}{Go for it}
\promptentry{23}{Final review}{2026-10-22 12:31 UTC}{What would we gain from fixing \#4 ?}
\promptentry{24}{Final review}{2026-10-22 12:40 UTC}{Let's fix it.}
\promptentry{25}{Final review}{2026-10-22 12:50 UTC}{Explain what we gain from fixing \#5 ?}
\promptentry{26}{Final review}{2026-10-22 12:56 UTC}{Can we fix this while staying as simple as we want to ?}
\promptentry{27}{Final review}{2026-10-22 12:56 UTC}{Go}
\promptentry{28}{Final review}{2026-10-22 12:59 UTC}{Excellent, now there is one issue I am concerned with, the perfomance of the model. What can we do to improve it ?}
\promptentry{29}{Final review}{2026-10-22 13:03 UTC}{Ok let's [Request interrupted by user]}
\promptentry{30}{Final review}{2026-10-22 13:03 UTC}{Let's go with A.}
\promptentry{31}{Final review}{2026-10-22 13:12 UTC}{How many lines of code/\allowbreak{}comments/Markdown are there in the full notebook ? \newline Make a short summary table.}
\promptentry{32}{Final review}{2026-10-22 13:17 UTC}{What are the results of the random forest model ? compare them to the tree}
\promptentry{33}{Final review}{2026-10-22 13:29 UTC}{Do we now have the final version of the notebook ?}
\promptentry{34}{Repository}{2026-10-22 13:38 UTC}{Yes let's now build the final README for this project, make it professional and visually beautiful. Do not make it too long.}
\promptentry{35}{Repository}{2026-10-22 13:40 UTC}{are these also clickable links ?\par ![Python](https:/\allowbreak{}/\allowbreak{}img.shields.io/\allowbreak{}badge/\allowbreak{}Python-3.12-3776AB?\allowbreak{}logo=\allowbreak{}python\&logoColor=white)\newline ![scikit-learn](https:/\allowbreak{}/\allowbreak{}img.shields.io/\allowbreak{}badge/\allowbreak{}scikit-{}-learn-1.9-F7931E?\allowbreak{}logo=\allowbreak{}scikitlearn\&logoColor=white)\newline ![pandas](https:/\allowbreak{}/\allowbreak{}img.shields.io/\allowbreak{}badge/\allowbreak{}pandas-3.0-150458?\allowbreak{}logo=\allowbreak{}pandas\&logoColor=white)\newline ![Data](https:/\allowbreak{}/\allowbreak{}img.shields.io/\allowbreak{}badge/\allowbreak{}data-UCI\%20·\%20CC\%20BY\%204.0-lightgrey)}
\promptentry{36}{Repository}{2026-10-22 13:41 UTC}{YES}
\promptentry{37}{Repository}{2026-10-22 13:43 UTC}{Let's push the folder to the following repo : git remote add origin https:/\allowbreak{}/\allowbreak{}github.com/\allowbreak{}yassineerraji/\allowbreak{}Augmented-Decision-Final-Project.git\par Before that, let's create a .gitignore to avoid tracking DS\_Store and .venv.}
\promptentry{38}{Slide deck}{2026-10-22 13:56 UTC}{Great, now we want to build a slide deck to present for our final presentation. \newline For context, it should not exceed 15 minutes and should include the essentials of our project.\par The names are Yassine ERRAJI, William PELTIER, Ben DAESH, Simon AMAR\newline We want the slides to be very clean and professional yet very original and not look like the main claude templates. It should stand out.\par It should be built in LaTeX and compiled in pdf, obviously following the classic landscape slide deck format.\par Include the very main formulas that we build the project on top of (avoid the too complex ones that do not appear in the lecture decks). \newline Include the most relevant visuals that we obtained from the notebook (make sure that its format is adapted to the slide deck).\newline Also include schemas or graphs that you create on your own if you believe it would help present better and help people understand better our project.\par One Important element, slides must be readable and not too dense, therefore, we would rather have more slides that are more readable and less content heavy that fewer dense ones.\par Let's do this inside of a slides folder at the root of the repo.}
\promptentry{39}{Slide deck}{2026-10-22 15:36 UTC}{In the slide 25/28, in the head of the plot, the text is not well displayed (our ratio: 10), fix it.}
\promptentry{40}{Report}{2026-10-22 15:43 UTC}{Okay great now we want to build a properly compiled PDF report containing:\newline •Methods and choices. Explain your approach. Motivate the model, preprocessing, and\newline hyperparameters, e.g., tree depth, leaf size, or regularization penalty.\newline •Evaluation. Compare with a baseline. Tune by validation or cross-validation, fitting\newline preprocessing and selection within each training split. Report final test performance.\newline •Evidence and recommendation. Include readable code snapshots, labeled plots, and\newline result tables that justify your final model or variable set, recommendation, and limitations.\par Build it in a folder called report}
\promptentry{41}{Slide deck}{2026-10-22 15:57 UTC}{Now I want you to reproduce another exactly same version of the slide deck we have but this time using the Arial font.}
\promptentry{42}{Slide deck}{2026-10-22 16:06 UTC}{Okay excellent, let's delete the previous version and keep the arial one as the main.}
\promptentry{43}{Appendix}{2026-10-22 16:37 UTC}{Great, we want to add an appendix to the report about LLM Assistance (take a look at the ProjectRequirements.pdf).\newline In that section we must explain exactly the procedure of the INSTRUCTIONS.md and AgentDiscussion folder. \newline Detail the workflow of plan and recap, my modifications and instructions etc..\par If possible, include the content of the files in a nice rendering.}
